\chapter{Conclusions and Recommendations}\label{ch:conclusions}

\section{Conclusion}\label{sec:conclusions}

This research set out to answer a practical question: can a natural-language
security request be carried end-to-end to an explainable, prioritised
vulnerability report with minimal human intervention, while keeping every step
inside a controlled safety boundary? The answer supported by the evidence is a
qualified yes. NeuroShell, as designed and evaluated across its four
components, covers every stage of the chain---intent recognition, planning,
execution and error recovery, and vulnerability analysis---and each stage returns
measured results on its own committed protocol.

The conclusions are therefore honest about their strength. The guarded intent
pipeline (C1) converts unguarded model output into a schema-enforced contract:
it raises decision accuracy by 8--13 points, converts a 0\% flaw-catch into a
72\%/64\% catch, and more than halves the false-positive rate under its own
1,500-case protocol. The planner (C2) demonstrates that a general instruction
model can be adapted to a structured intent-to-command task and grounded in
maintainable documentation: 87.54\% strict exact-match on the held-out set, and
a 28-point semantic-accuracy gain from retrieval in the separate base-model
ablation. The error-recovery engine (C3) establishes the framework's
distinctive capability: reading a failed security tool's error, labelling it,
and selecting a fix on the device, with a 99.1\% macro-F1 at a realistic 30\%
message-noise level. The vulnerability analysis engine (C4) shows that a
compact, explainable classifier over five vulnerability-intelligence features
can separate vulnerabilities with documented exploitation evidence, subject to
the provisional figures being confirmed (93.00\% accuracy, 87.47\% F1, 97.20\%
ROC-AUC, 90.30\% PR-AUC).

The honesty of the reporting is itself a conclusion. Every measured figure is
reported with its verification context; targets that were not measured (the
AEERE fix-rate, speed, and safety-tables; the AVAE confirmed final-model
evaluation; an integrated end-to-end benchmark) are left open in the component
reports and in this dissertation rather than filled with estimates. This means
the current strongest claim is that the framework behaves as designed against
its own committed, reproducible evaluations, not that it has been validated
against an independent benchmark or by end users.

\section{Limitations}\label{sec:conclusions-limitations}

Six limitations bound the conclusions. (1) The evaluation data sets (IRE
1,500-case set, planner 337-case held-out set, AEERE 8,000-case set, AVAE
155,852-record set) were authored or curated by the research group, so the
reported rates measure performance on the components' own labelled
distributions. (2) The results are single-run point estimates on one CPU-only
machine; no cross-platform or cross-hardware generalisation is established.
(3) No component completed a formal user or stakeholder validation study, so
usability and deployment-practice claims remain untested. (4) The planner's two
experiments measure different configurations and cannot be combined into a
single end-to-end accuracy figure. (5) The AEERE fix-rate, end-to-end speed, and
dangerous-command block-rate targets could not be measured because the live
engine (Docker plus Ollama) was not available at the time of writing. (6) The
AVAE performance figures are provisional and depend on confirmation of the final
model configuration.

\section{Recommendations}\label{sec:conclusions-recommendations}

The following recommendations follow directly from the findings and
limitations.

\begin{enumerate}
  \item \textbf{Close the open engineering gaps before deployment.} The IRE's
  default development API key and inactive \emph{production} scope mode, and the
  planner's pending deprecation warning, are the concrete items identified so
  far; a pre-deployment security review should look for further ones.
  \item \textbf{Measure the live AEERE engine.} Execute the committed SO1, SO2
  and SO4 target runs on the full engine (Docker plus Ollama) and complete the
  fix-rate, speed, and dangerous-command-block tables that the component report
  intentionally left blank.
  \item \textbf{Confirm the AVAE model results.} Complete the verified final
  evaluation of the XGBoost configuration and replace the provisional metric
  table with the confirmed one, keeping the same data description.
  \item \textbf{Conduct an integrated evaluation.} Design a factorial comparison
  (the same unseen cases across the pipeline stages) so that adaptation,
  retrieval, execution recovery, and analysis can be separated, and so that a
  single end-to-end benchmark for the whole framework can be reported.
  \item \textbf{Build an independent test corpus and user study.} An
  independently authored prompt/error corpus and a formal user evaluation of the
  CLI and API would lift the current self-authored evaluation limitation.
  \item \textbf{Maintain the knowledge base as a first-class artefact.} Because
  the planner keeps documentation separate from the model, updates to tool
  syntax do not require retraining; a review process for the ChromaDB store
  should be defined as part of maintenance.
\end{enumerate}

\section{Future Work}\label{sec:conclusions-future}

Each component names its own next steps, and the framework as a whole inherits
them. The IRE should close the API-key and scope-mode gaps IRE names, extend
the data set to additional tools and intent classes, and run a user study. The
planner should extend RAG evaluation to the final adapted model on a factorial
design, broaden the knowledge base to more tools and option patterns, and move
from a strict exact-match criterion toward a documented semantic evaluation that
distinguishes valid command alternatives from errors. The AEERE should complete
the live-engine measurements, extend the error list and tool-swap rules from
logged UNKNOWN cases as real tools change, and test recovery under deliberately
adversarial error messages beyond the current noise injection. The AVAE should
confirm its final model, add broader scanning-tool input support, and expand the
data set with additional exploitation-related signals.

The framework-level next steps are to run the components as one integrated
system, to test the hand-over contracts (intent contract, planner output,
execution result, analysis output) under a single end-to-end scenario, and to
use the SQLite safety logs and Neo4j findings graph as the audit trail for a
realistic depth-first engagement against a safe lab target. If those runs
sustain the component-level results, the framework's core claim---an autonomous,
local, guarded natural-language-to-vulnerability-report chain---will be
supported end to end, and the open gaps identified in this dissertation will be
closed.

This dissertation set out to remove the human from the loop where it is safe to
do so---between intent and command, between command and execution, and between
execution and report---while leaving the human in the loop of what to engage.
That is the boundary NeuroShell is designed to respect, and the results and
limitations reported here delineate how far, exactly, it currently goes.